% !TEX root = ../main.tex
\chapter{Formas de onda: GTKWave e Surfer}
\label{cap:ondas}

Depois que a simulação gera o arquivo de ondas, a AURORA abre um visualizador para inspecionar cada sinal do processador ciclo a ciclo. Há dois visualizadores disponíveis. O \tool{GTKWave} é o padrão e já vem empacotado: um \textit{fork} próprio do laboratório, o \repo{gtkwave-nipscern}, com tema escuro e ajustes de usabilidade. O \tool{Surfer} é um visualizador moderno escrito em Rust, opcional, do qual a AURORA usa o \textit{fork} \texttt{surfer-aurora}, mantido pelo grupo no GitLab com decodificadores específicos do SAPHO. Ambos leem o mesmo formato de despejo, o FST, de modo que trocar de visualizador não exige simular de novo.

\figurapendente{Formas de onda de uma simulação do processador media\_movel no GTKWave}{Mostrar as faixas de clock, as trilhas de Assembly e C± em sincronia, as variáveis e as portas.}

\section{O que a AURORA prepara para você}
\label{sec:curadoria-ondas}

Abrir um despejo bruto num visualizador qualquer significa começar do zero, sem nenhum sinal selecionado e tudo em binário. A AURORA elimina esse trabalho gerando, a cada simulação, um \textit{layout} curado. Cada processador do projeto vira uma seção própria, com divisores nomeados e, no Surfer, um grupo colapsável. Os sinais de topo, \textit{clock}, \textit{reset} e interrupção, vêm primeiro; seguem as portas de entrada e saída em amarelo, as instruções em violeta, as variáveis em laranja e as \textit{flags} ao final. Variáveis do tipo \code{comp} são decodificadas para a forma legível $a+bi$ com o conversor \tool{comp2gtkw} do YANC.

O destaque são as trilhas de texto que mostram, a cada ciclo de \textit{clock}, o mnemônico \textit{assembly} executado e a linha \cmm{} correspondente, avançando em sincronia com o circuito: é o seu programa rodando dentro do hardware, visível linha a linha. Essas trilhas vêm dos arquivos de tradução gerados na compilação. No GTKWave, toda essa curadoria chega como um arquivo \file{.gtkw}; no Surfer, como um arquivo de estado \file{.surf.ron}. Nos dois casos, o visualizador já abre pronto para leitura.

\section{Escolhendo o visualizador}

O seletor de visualizador fica na \textit{toolbar}, ao lado do seletor de simulador, e vale para o aplicativo inteiro, lembrado entre sessões. Se o Surfer estiver selecionado mas o executável não estiver presente, a AURORA recua para o GTKWave sem erro.

\section{GTKWave}
\label{sec:gtkwave}

O GTKWave empacotado não é o original puro. O \textit{fork} do laboratório traz tema escuro por padrão, com o \textit{canvas} na paleta da AURORA, ajuste automático de \textit{zoom} ao abrir, de modo que a onda inteira caiba na tela, rótulos justificados à esquerda, ícones modernizados e \textit{zoom} animado.

No uso cotidiano, os sinais são adicionados a partir da árvore de escopos à esquerda; o cursor primário se posiciona com o clique e o secundário com o botão do meio, com a barra exibindo a diferença de tempo entre eles; o formato de exibição de cada sinal muda pelo menu de contexto, em \menu{Data Format}; o \textit{zoom} responde aos botões de lupa e a \tecla{Ctrl} com a roda do mouse; e a seleção atual de sinais, cores e ordem pode ser salva em um \file{.gtkw} próprio por \menu{File}, \menu{Write Save File}. Após uma nova simulação, \menu{File}, \menu{Reload Waveform} recarrega o arquivo.

\begin{nota}
Se você salvar o seu próprio \file{.gtkw} no projeto e o marcar como ativo no seletor da AURORA, ele passa a ter precedência sobre o \textit{layout} automático: a simulação seguinte grava exatamente os sinais que o seu arquivo referencia (Capítulo~\ref{cap:simulacao}).
\end{nota}

\section{Surfer}
\label{sec:surfer}

O Surfer é um visualizador com interface acelerada por GPU e carregamento preguiçoso dos arquivos, confortável em simulações grandes. A AURORA baixa o binário do \textit{fork} \texttt{surfer-aurora} automaticamente, com verificação de integridade, e o lança com o \textit{layout} curado.

Entre os recursos que valem conhecer estão os formatos de exibição ricos, de binário e hexadecimal a ponto fixo e ponto flutuante IEEE em vários tamanhos; os grupos colapsáveis, um por processador no \textit{layout} da AURORA, essenciais em projetos multiprocessados; as trilhas de \textit{assembly} e \cmm{}, instaladas automaticamente como tradutores de mapeamento; os cursores e a navegação rápidos, com uma linha de comando interna de autocompletar difuso; e o desenho analógico de sinais, útil para ver amostras de sinais processados como curvas. Os complexos do SAPHO também funcionam: a AURORA pré-decodifica os sinais \code{comp} e instala o resultado como tradutor, de modo que a faixa já abre na forma $a+bi$.

Uma limitação conhecida da versão empacotada no Windows é o recarregamento automático do arquivo de ondas, que não dispara; por isso a AURORA fecha e reabre a janela do Surfer a cada nova simulação. No modal de configuração de ondas há uma opção de manter as janelas abertas, para comparar execuções lado a lado.
